\section{Práctica de ingeniería de agentes: planeación, ejecución y verificación}
\subsection{Detalles de Plan-Act-Verify}
En la ingeniería real, el agente debe contar con un ciclo de acción explícito: primero leer el issue, invocar el servicio de recuperación para acotar el contexto, planear el parche, ejecutar la modificación del código, correr las pruebas y luego ajustar con la retroalimentación de los registros. Graham enfatiza la combinación de varias herramientas: el LLM se enfoca en generar, los scripts se encargan de operar, el framework de pruebas verifica los resultados y, al final, la evidencia se sube al dashboard para la revisión humana.

Señala que, en cada ronda de Plan, primero debe generarse una lista de tareas que indique explícitamente qué archivos se van a modificar y qué pruebas se van a correr. En la etapa de Act se vinculan scripts y comandos de terminal concretos para que el agente escriba y ejecute el código de verdad, y en la etapa de Verify se corren las pruebas y el análisis estático; las muestras que fallan disparan una nueva Plan.

\subsection{Colaboración humano-agente y cobertura de herramientas}
\begin{figure}[H]
\centering
\includegraphics[width=\textwidth]{frame-03-human-loop.jpg}
\caption{El supervisor humano puede pausar la ejecución del agente, revisar sus decisiones e intervenir, lo que permite corregir errores en tiempo real. \protect\footnotemark}
\end{figure}
\footnotetext{Intervalo de tiempo en el video: 00:32:05--00:32:10.}

OpenHands permite que el desarrollador inserte una revisión humana en cualquier paso: cuando el agente planea ejecutar un comando peligroso (como un despliegue o una escritura en la base de datos), el flujo puede pausarse, revisar los registros y, de ser necesario, modificar el prompt o ejecutar manualmente parte de los comandos. Las tareas de larga duración también envían checkpoints, lo que facilita dar seguimiento al estado.

\begin{importantbox}{Mejor práctica}
Implementa un pipeline para el agente: primero plan (segmentar la tarea, encontrar las dependencias), luego act (edición de código, compilación, pruebas) y por último verify (registros, instantáneas, capacidad de reversión). Los pasos que fallan necesitan retroceder y transmitir el contexto clave al LLM.
\end{importantbox}

\begin{knowledgebox}{El valor de la intervención humana}
Permitir que un humano observe el comportamiento del agente durante su ejecución, y que pause antes de que se realice una operación peligrosa, es clave para mejorar la confiabilidad, sobre todo en operaciones de compilación de lanzamientos o de base de datos.
\end{knowledgebox}

\begin{warningbox}{Ventana de contexto y alucinaciones}
El contexto del agente es limitado; si se envía de una sola vez toda la estructura del repositorio, se rebasa la ventana legible, así que primero debe hacerse extracción o recuperación de información, pues de lo contrario el LLM producirá parches alucinados por falta de información del entorno.
\end{warningbox}

\subsection{Resumen de la sección}
\begin{itemize}
    \item El ciclo de planeación, ejecución y verificación hace que el agente sea controlable dentro del proceso de ingeniería.
    \item La colaboración humano-agente (pausar, revisar registros, saltar al contexto) es un medio necesario para evitar errores graves.
\end{itemize}
